\section{A completely specified coefficient instance}
\label{sec:coefficient-instance}

The coefficients in this construction are fixed integers, not existential
parameters. Their binary expansions are too long to print or materialize.
This section defines them by finite sums and fixed integer powers. The
attached standard-library programs implement precisely the same definition.
Every loop in coefficient construction is bounded before the computation
starts; no test of whether K215 halts enters the definition.

\subsection{The literal machine and its one-sided normalization}

We pin ProveIt at commit
\texttt{cdd582bec13f908df3711fa9e05d1ddac368a5db}.
The delivered native table has 215 working states, start 0 and halt 254.
Its dense presentation has start 0, working states 0 through 214, and halt
215. The supplied \texttt{canonicalization.json} is checked against all
430 transition rows in both directions; its two table hashes are checked
against literal expected hashes. Thus normalization begins with a state
renaming of the exact dense machine, not with a related candidate.

Fold the original tape onto the nonpositive physical coordinates. At
physical position $n\leq0$, the two stored bits represent original
positions $-n$ and $n-1$, respectively. Four ordinary symbols represent
the bit pairs $00,01,10,11$; four marked versions represent the same pairs
at the origin. Ordinary $00$ is blank symbol 0. Add input symbol 1 and
temporary origin flag 2, giving ten tape symbols in total.
For each native working state $q$ and track $s\in\{0,1\}$, introduce a
state $(q,s)$. On track 0 an original right move becomes a physical left
move, and an original left move becomes a physical right move except at
the marked origin, where it stays and changes to track 1. On track 1
these directions are interchanged; an original right move at the marked
origin stays and changes to track 0. Only the selected bit is overwritten.
An original transition into 254 enters the single normalized halt state.

There are six additional states: $q_0,q_1$, erase, return, halt and trap.
The first instruction on input 1 writes flag 2 and \emph{stays} in the
distinct nonhalting state $q_1$. It then moves left, erases every remaining
initial 1, reverses at the first blank, returns to flag 2, replaces it by
marked $00$, and starts the folded native state $(0,0)$. Halt and trap are
absorbing stay states; unspecified malformed cases enter trap.
On an initial block of $n\geq1$ ones ending at the origin, this prelude
takes exactly $2n+2$ instructions and leaves a completely blank folded
tape with one marked origin. It never visits a positive coordinate.

The resulting machine has 436 states and 4360 total instructions. A
direct induction on source steps proves its halting equivalence with
blank K215, for every input-block length. The finite checker additionally
checks every one of the 3440 combinations of native state, track, pair
and origin status, as well as 32 prelude lengths. The machine need not
reject odd input lengths: our fixed arithmetic instance selects four
initial ones, so behavior on the other lengths is immaterial.

\subsection{The fixed window alphabet}

Use the helical tableau clauses from the repository's
\texttt{EXPLORATION\_HELICAL\_UNARY\_TABLEAU.md}, with the subsequent
stay-step markers of \texttt{FIXED\_RAW\_UNIVERSAL\_81\_PROOF.md}.
There are the two boundary tiles $V,H$; for every tape symbol and each
head state or no head, there is an ordinary unphased tile; finally there
are the four initial phase tiles $L,I,Q,R$. Order these respectively as
$V,H$, then symbol first and head second (no head first), then $L,I,Q,R$.
This is an explicitly ordered alphabet of size
\[
        a=2+10(436+1)+4=4376.
\]

The complete local predicate is given by \texttt{local\_valid} in
\texttt{fixed\_machine.py}. Its clauses propagate vertical boundaries,
propagate horizontal boundaries only between non-$V$ cells, disallow
adjacent boundaries, enforce initialization $L^rI^uQR^s$, enforce the
exact three-neighbor head update, forbid a head escaping a vertical
boundary, and require a halted head in the last interior row. These are
exactly finite radius-one conditions. The three-by-three window's four
corner entries are not tested by its center predicate; they are retained
and identified by the exact overlap conditions of the lifted system.

Let $S$ be the Start window
\[
\begin{array}{ccc}
H&H&H\\
(1,\varnothing,I)&(1,\varnothing,I)&(1,q_0,Q)\\
(1,\varnothing)&(1,\varnothing)&(2,q_1)
\end{array}
\]
and let $E_0$ be the End window
\[
\begin{array}{ccc}
H&H&H\\
(0,\varnothing,L)&(1,\varnothing,I)&(1,\varnothing,I)\\
(0,\varnothing)&(1,\varnothing)&(1,\varnothing).
\end{array}
\]
Both satisfy the actual predicate. Enumerate admissible windows by
$W_0=S$, $W_1=E_0$, followed by all remaining admissible windows in the
lexicographic order of their nine tile indices. Define $k$ to be their
number. Factoring the five-cell cross predicate by the incoming-head
histograms gives
\[
 \#\{\hbox{admissible crosses}\}=314330772961757,
 \qquad
 k=115264654755851622862362693632.
\]
The latter is the former multiplied by $a^4$, for the free corners.
The attached counter is also compared with direct enumeration on four
small machines, covering 995328 cross assignments. The mathematical
definition of $k$ remains the displayed finite window enumeration.

The repository's marked helical tableau argument applies to this
alphabet. A cyclic word with exactly one Start at position 0, an End at
position 2, and the exact lifted relation at offsets $0,-1,-h$ decodes a
finite halting run from $L^rI^3QR^s$. Thus there are four initial ones.
Conversely any such halting run has a padded helical presentation with
independently enlargeable spatial width and time height. This is the
precise semantic input needed below.

\subsection{The exact coefficient formulas}

Choose the clause radix $A=2^{100}$ and write $\ell=100$.
In the formulas below $r,c\in\{0,1,2\}$ and $s\in\{0,\ldots,a-1\}$.
Define
\[
 j(r,c,s)=1+(3r+c)a+s,\qquad
 p(r,c,s)=k+3a+(2-r)3a+(2-c)a+s.
\]
Put
\begin{align*}
 J_0&=9a+4,&m_0&=2\ell+30a+8,\\
 p_0&=\max\!\left(0,\left\lceil\frac{k+9a+4-m_0}{2}\right\rceil\right),
 &J_*&=J_0+p_0,\\
 m&=m_0+2p_0,&K&=m+1.
\end{align*}
For this instance $K=k+9a+5$, so exactly one ignored dummy is used.
Its position is $e_*=k+12a$. Set
\[
 M=k+15a+1,\quad
 (v_1,v_2,v_3,v_4)=(M,3M,9M,27M),\quad E_{max}=27M,
\]
and let $\mathcal E$ be the disjoint union of selector positions
$0,\ldots,k-1$, all $p(r,c,s)$, $e_*$, and the four $v_i$.
The nonzero center-clause coefficients are
\begin{align*}
 c_i&=1+\sum_{r,c}A^{1+(3r+c)a+W_i[3r+c]}
               +\sum_{l=1}^4A^{9a+l} &&(0\leq i<k),\\
 c_{p(r,c,s)}&=A^{j(r,c,s)},&
 c_{v_l}&=A^{9a+l} &&(1\leq l\leq4).
\end{align*}
Every other $c_e$, in particular $c_{e_*}$, is zero. Define
\[
 \mu=A-2+\sum_{j=1}^{J_*}A^j,
 \quad H=51M+3a+1,
 \quad T_1=H+2E_{max}+a+1,
 \quad T_2=T_1+2E_{max}+1.
\]
Take
\[
 b=L=5^{45},\qquad d=bL=5^{90},\qquad
 R=2^b,\qquad B=R^L=2^d.
\]
The exact remaining integers are
\begin{align*}
 D_C={}&R^{3a}+R^{H+a}+R^{8M}+R^{24M}
       +\sum_e c_e\bigl(R^{T_1-e}+R^{T_2-e}\bigr),\\
 D_R={}&R^H,\\
 M_C={}&B-1-\sum_{e\in\mathcal E\setminus\{1\}}R^e-2R^{e_*},\\
 M_F^{\rm nat}={}&\mu(R^{T_1}+R^{T_2})
       +\sum_{r=0}^1\sum_{c=0}^2\sum_{s=0}^{a-1}R^{p(r,c,s)}\\
       &+\sum_{r=0}^2\sum_{c=0}^1\sum_{s=0}^{a-1}R^{H+p(r,c,s)}
       +R^{9M}+R^{27M}+4.
\end{align*}
Thus the fixed ports used by the specialized polynomial can be set to
\[
 B-1,\quad D_C+BD_R,\quad M_C,\quad M_F^{\rm nat}+B-1,
 \quad W_0=2^{2\cdot5^{90}+5^{45}},\quad W_0+2\cdot5^{90}.
\]
All sums have specified finite ranges, and every exponent is a fixed
integer. None is an exponentiation of an unknown in the resulting
Diophantine equation.

\subsection{Why these are admissible compiler parameters}

The repository programs choose small sufficient clause and native
radices; their mathematical proofs only require the stated inequalities.
Here $A>k+1$, so every selector sum and selector-plus-copy clause has
radix-$A$ digits below $A$. The low digit $A-2$ and the remaining unit
digits of $\mu$ consequently implement exactly the same clause tests.
Moreover
\[
 \operatorname{popcount}(\mu)=\ell-1+J_*,\qquad
 m=2\operatorname{popcount}(\mu)+12a+2,
 \qquad |\mathcal E|=m+1.
\]
Padding clauses have identically zero expressions and create no extra
condition on actual windows. The one extra dummy has zero coefficient
in every tested field.

For an explicit sufficient radix bound put $n_c=k+9a+4$ and
$G=128(K+2)n_c+16$. Each $c_e\leq14A^{J_*}$, hence
$\sum_e c_e\leq14n_cA^{J_*}$, while $\mu<A^{J_*+1}$.
The delivered numerical geometry checks
\[
 b\geq\ell(J_*+2)+\operatorname{bitlength}(G).
\]
It follows that $R>A^{J_*+2}G$ and therefore
\[
 R\geq\max\{4(K+2)(2\sum_ec_e+6)+8,\;2\mu+4,\;16\}.
\]
This is the stronger half-binomial compiler margin, including the
possible extra upper dummy bit. All unshifted raw coefficients are at
most $R/4-2$; a rotated native digit is at most $3R/4$, or $R/2+1$ at
the spill residue. Thus no carry is hidden by the chosen larger radix.

Set $g=T_2+E_{max}+1$. The explicit value $L=5^{45}$ satisfies
$L>g+E_{max}$. The two center-clause bands are separated by
$2E_{max}+1$, and their complementary cyclic gap exceeds $E_{max}$.
The four anchors retain the unique difference $18M$, and
$e_*+1<E_{max}$, so the original synchronization proof, including the
half-binomial dummy-spill case, applies verbatim to these positions.

Here $A\equiv1\pmod5$, so every selector coefficient is $14$ modulo
5 and every payload or anchor coefficient is 1 modulo 5. Both $k$ and
$a$ are divisible by 4. The contiguous selector and payload intervals
therefore vanish in each weighted sum modulo 5. Since $M\equiv1\pmod4$,
the anchor exponents have residues $1,3,1,3$, which also cancel.
The four low shifts in $D_C$ are each 0 modulo 4. Because $b$ is a
power of 5, $R\equiv2\pmod5$, and hence
\[
      D_C\equiv4,\qquad D_R\equiv1,\qquad
      2D_C-D_R\equiv2\pmod5.
\]
No optional high correction monomial is needed. This verifies the
five-adic unit used to adjust the actual packed index.

Finally $0<M_C,M_F^{\rm nat}<B-1$, $M_C$ is even, and
$M_F^{\rm nat}\equiv4\pmod8$. The baseline low mask has $d-m$ set bits;
clearing the second dummy bit decreases that by one. The baseline high
mask has $m$ set bits; adding its previously absent bit 2 increases
that by one. Consequently
\[
       \operatorname{popcount}(M_C)+
       \operatorname{popcount}(M_F^{\rm nat})=d.
\]
The source coefficient is the shifted $M_F^{\rm nat}+B-1$, as required
by the complete75 and subsequent polynomial interfaces. Treating that
shifted coefficient itself as a one-cell mask would be incorrect.

\subsection{Reproducibility and limits of the checks}

\texttt{fixed\_machine.py} reproduces the normalized table and machine
receipt. \texttt{coefficient\_recipe.py} reproduces the compact exact
coefficient specification. \texttt{verify\_bridge.py} reruns their finite
checks, the direct cross-counter comparisons, and one fully materialized
small compiler example. The function \texttt{materialize\_k215} gives a
literal terminating coefficient algorithm but is intentionally not run:
already $B=2^{5^{90}}$ would require about $5^{90}$ bits.
These checks establish finite interface facts and corroborate the
mathematical compiler proof. They are not a proof-assistant check of the
entire Diophantine equivalence, and no halting run or full Pell witness
for K215 is asserted or numerically materialized.
